\subsection{Background}\label{subsec:background}

The JPEG format is a widely used and well-understood standard for compressed image storage. 
In contrast, the PPM format is less commonly used in practice but is highly relevant in experimental contexts.
PPM images store uncompressed pixel data in a simple and explicit representation, allowing direct access to color values without the need for complex decoding.

These properties make PPM a suitable format for studying low-level image processing algorithms, as the distinct transformation of each pixel is possible.
Therefore, understanding the structure and limitations of the PPM format is essential for designing correct and efficient denoising and binarization algorithms.
